\section{The fiscal condition}
\label{sec:fiscal}

If income moves from wages to AI capital, the wage taxes on it stop being collected and whatever
tax reaches the capital is collected instead. The fiscal condition asks whether the second covers
the first. This section builds both sides from components and reports the answer under every
reading we could defend, because the answer turns out to depend on choices that reasonable people
make differently and we would rather show the whole surface than pick a point on it.

\subsection{What the condition requires}

Not every dollar of displaced wages is lost to the wage tax base. Some of the displaced are
re-employed, at a wage that is on average lower than the one they lost, and what they earn is
taxed as before. We call the share of the displaced wage bill that comes back as wages the
\emph{retained wage share}; measured on fourteen biennial vintages of displaced worker data and
solved against the stock-flow identity described in Section~\ref{sec:data}, it is
\result{RetainedWageShare}, and it is replicated.

The condition then has a simple form. Writing $\tau_\ell$ for the effective tax rate on wages,
$\tau_k$ for the effective rate on the capital income that replaces them, and $R$ for the retained
wage share, replacing the lost revenue requires
\begin{equation}
\tau_k \;\ge\; \tau_\ell \,(1 - R).
\label{eq:condition}
\end{equation}
The right-hand side depends on which effective labor tax rate is used, and two defensible readings
exist, a narrower one from the automation tax literature and a broader bottom-up one built from
federal taxes over wages and salaries. Across them the required rate runs from
\result{RequiredEasy} to \result{RequiredHard} percent. We carry both throughout and call the lower
one the easier reading, because it is the one the rate has the best chance of clearing. This much
is replicated: an independent rebuild reproduced the retained wage share to six decimal places and
both bounds of the required rate.

\subsection{Building the rate on the income that shifts}

The quantity on the left of Equation~\eqref{eq:condition} is not a statutory rate and not an
economy-wide average. It is the rate actually borne by a marginal dollar of surplus generated by
AI capital in the United States, counting every layer of tax exactly once: the corporate tax at
the entity level, federal and state, net of what is shifted to lower-taxed jurisdictions, and then
the tax that the owner of the claim pays on what reaches them.

Two features of that build do most of the work, and neither is our assumption.

The first is expensing. Under full expensing of equipment the normal return to an investment is
already relieved at the entity level, so at the margin it bears only the tax paid by the owner.
That is not a modelling convenience; it is the result in \citet{AcemogluManeraRestrepo2020}, and we
take it from them rather than deriving it again. For an equity-financed corporation the normal
return therefore bears the shareholder rate alone. For a debt-financed one it bears the rate on the
bondholder less the corporate rate, which at sourced values is negative throughout. A negative rate
on the normal return is not an error: it is what the interaction of expensing and interest
deductibility produces, and the Congressional Budget Office's own measured effective rate on
debt-financed corporate investment, minus \result{CBODebtFinancedETR} percent, sits inside the
interval our components imply.
The sign was predicted by the theory and the magnitude measured by the budget office, and our
assembly reproduces both from components.

The second is the rent share. Only the part of the return that is economic rent bears the full
entity-level rate, and the literature does not agree on how large that part is. We carry two
readings and never average them, because they are not two estimates of one quantity but two
incompatible readings of what has happened to the corporate sector. The first puts the rent share
at \result{RentBarkai}; the second puts it at zero, in which case the assembled rate collapses to
the rate on the normal return alone. Reporting a midpoint would invent a position neither
literature holds.

\paragraph{The base, and the argument against our own choice.} We take the rate on the marginal
flow of AI surplus, whoever ends up holding the claim on it. The case against is worth stating in
full, since it is the strongest objection to this section. An economy-wide average effective rate
on capital income, of the kind the International Monetary Fund reports for advanced economies at
twenty to twenty-two percent, is measured on a real capital stock that is mostly neither fully
expensed nor easily shifted, and includes personal-level taxes on dividends and realized gains as
they are actually paid. If one believes AI capital will come to look like the rest of the capital
stock, that is the relevant rate, and the condition closes on it. We think the marginal flow is
the right object for a revenue replacement question, because the question is what happens to the
next dollar that moves, not what the average existing dollar bears. But this is a judgement about
which base answers which question, not a fact, and a reader who takes the other view should read
our result as bounding the pessimistic case rather than settling the matter.

\subsection{The answer, under both readings and both bases}

Table~\ref{tab:tauk} gives the assembled rate against the required band. At the centre, under the
first rent reading, the rate on income moving from wages to AI capital is \result{TauKBarkai}
percent, against the \result{RequiredEasy} to \result{RequiredHard} percent required. Under the
second reading it is \result{TauKKN} percent. Both fall short, and they fall short under both
readings of the labor tax rate. This is a scenario result in the strict sense: three of its
components are swept rather than sourced, and the sweep is what the table reports. The
construction itself is replicated: every sealed quantity for it matched in round three, inside the
tolerance set before the round began.

\begin{table}[htbp]
\centering
\caption{The assembled effective rate on income that moves from wages to AI capital, against the
rate required to replace the wage taxes lost. Scenario, on a replicated construction.}
\label{tab:tauk}
\begin{tabular}{llrrrr}
\toprule
Base for the taxable holder share & Rent reading & Central & Fifth to ninety-fifth & Analytic maximum & Share of the space that closes \\
\midrule
All equity outstanding & First rent reading & 0.0851 & 0.0726 to 0.0993 & 0.1236 & 0.0013 \\
All equity outstanding & Second rent reading & 0.0150 & 0.0040 to 0.0255 & 0.0403 & 0.0000 \\
Domestically held stock & First rent reading & 0.1024 & 0.0860 to 0.1231 & 0.1506 & 0.2832 \\
Domestically held stock & Second rent reading & 0.0322 & 0.0165 to 0.0507 & 0.0685 & 0.0000 \\
\midrule
Required, easier labor tax reading & & 0.1101 & & & \\
Required, harder labor tax reading & & 0.1373 & & & \\
\bottomrule
\end{tabular}

\end{table}

How far short depends on the base for the share of corporate equity held in taxable accounts, and
the two answers differ enough that we report both and say which question each answers. On our base,
that share is measured over all US equity outstanding, foreign holders included, because a foreign
holder bears close to no US tax at the owner level and the dollar they hold is still part of the
surplus. On that base no corner of the parameter space closes the condition: the analytic maximum
over every combination of the swept components is \result{TauKAnalyticMax} percent, below even the
easier threshold. On the alternative base, which is the one the Congressional Research Service uses
and which measures the same numerator over domestically held stock alone, the share rises to
\result{CRSTheta} percent, the central rate rises to \result{TauKCRSCentral} percent, and about a
quarter of the parameter space, \result{TauKCRSPassShare} percent of it, closes the condition
against the easier labor tax reading. The centre still falls short on either base and under either
rent reading, and under the second rent reading the condition fails everywhere on both bases.

The difference between the two bases is not a disagreement about data. Both use the same numerator
from the same source. They differ over whether the foreign-held slice of equity belongs in the
denominator of a question about how much tax reaches the owner. For a question about
the average domestic shareholder it does not. For a question about how much tax a dollar of US
surplus bears, it does, because that dollar exists whoever owns the claim. That is why our
statement about the parameter space carries its base with it wherever it appears, and why it is
recorded as a limitation rather than a footnote.

\subsection{Why the rate is low, and it is not profit shifting}

The interesting part of this section is not the number but the reason for it. The binding
constraint is who owns the claims, not where the profits are booked.

About \result{ThetaUntaxable} percent of US corporate equity is held in forms the individual income
tax barely reaches: pension funds and retirement accounts, foreign investors, and nonprofits. Only
\result{ThetaTaxable} percent sits in taxable accounts, across a sourced range of
\result{ThetaTaxableLow} to \result{ThetaTaxableHigh} percent. Of the gains that do accrue in
taxable hands, a large share is never taxed at all: \result{GainsUntaxedAtDeath} percent of accrued
gains are held until the owner dies, at which point the basis is stepped up and the gain
disappears from the tax base permanently. What is realized is realized late, and the deferral
between accrual and realization cuts the present value of the tax by a factor that runs from
\result{DeferralLow} to \result{DeferralHigh} with a central value of \result{DeferralCentral}. The
same structure holds on the debt side: \result{InterestSheltered} percent of corporate interest is
received by holders who do not pay tax on it, chiefly retirement plans, and
\result{BondsRestOfWorld} percent of corporate bonds are held by the rest of the world, against
\result{BondsHouseholdDirect} percent held directly by households.

Put together, these are why the tax on the owner's side adds only
\result{ShareholderLayerPct} percent to the rate on rents rather than the figure that headline
dividend and capital gains rates would suggest.

Profit shifting is real and it is not the explanation. At the sourced \result{ShiftedShare} percent
of rents booked in lower-taxed jurisdictions, moving that share across its whole range changes the
assembled rate by \result{SwingShifted}, less than the range of the deferral factor at
\result{SwingDeferral} and less than the state corporate tax at \result{SwingStateCIT}.
Table~\ref{tab:levers} ranks every lever by how much it actually moves the rate. Profit shifting is
fourth of seven, and the levers above it have had a small fraction of the attention.

\begin{table}[htbp]
\centering
\caption{The levers, ranked by measured influence on the assembled rate. Each is moved across its
whole sourced or swept range with everything else at its central value. Scenario.}
\label{tab:levers}
\begin{tabular}{lrrr}
\toprule
Lever & Rate at the low end & Rate at the high end & Swing \\
\midrule
Deferral and step-up at death & 0.0772 & 0.0931 & 0.0159 \\
State corporate income tax & 0.0785 & 0.0918 & 0.0133 \\
Profit shifting & 0.0926 & 0.0802 & 0.0124 \\
Shareholder-level statutory rate & 0.0794 & 0.0909 & 0.0115 \\
Debt share of AI capital & 0.0883 & 0.0820 & 0.0063 \\
Effective rate on bondholders & 0.0830 & 0.0872 & 0.0042 \\
Share of equity in taxable accounts & 0.0823 & 0.0861 & 0.0037 \\
\bottomrule
\end{tabular}

\end{table}

One trap in that table deserves a sentence, because it is easy to read backwards. The share of
equity in taxable accounts appears last, with a swing of \result{SwingTheta}. That is because it
has been pinned to a narrow sourced range, not because it does not matter. It is simultaneously the
largest single determinant of the level of the rate and the smallest remaining source of
uncertainty about it. Sourcing a parameter removes it from the variance decomposition without
removing it from the mechanism.

\subsection{Two published figures, checked}

A rate assembled from components should reproduce something somebody else has published, and ours
reproduces two.

The first is at the owner level. Our shareholder layer at central values is
\result{ShareholderLayerPct} percent, and the Congressional Research Service's own text puts the
overall effective capital gains tax rate on corporate profits at about \result{CRSTextFigure}
percent. That is a direct corroboration and it stands. The same report's published table gives a
band of \result{CRSBandLow} to \result{CRSBandHigh} percent, and our layer does not reproduce it as
printed. It should not: that band already has the service's own taxable-share adjustment applied
over domestically held stock, while ours applies a share measured over all equity outstanding.
Rescaled to a common base the band is \result{CRSRescaledLow} to \result{CRSRescaledHigh} percent
and our \result{ShareholderLayerPct} percent sits inside it. We report the correction rather than
the original comparison, because as first written our own check failed against our own construction
and the failure was ours.

The second is on the debt side and is described above: the interval our components imply for the
debt-financed normal return contains the budget office's measured rate.

\subsection{What this section does not reach}

Foreign holders are treated as bearing no tax at the owner level. In fact dividends paid to them
bear US withholding tax at treaty rates, which are lower than domestic rates but not zero. The
effect on the assembled rate is upward and we have not quantified it. Since the result is that the
rate falls short, an unquantified upward correction is a limitation that runs against our own
finding, which is why it is stated here rather than left to the appendix. It is recorded with the
other limitations in Section~\ref{sec:limitations}.

Two further things are outside this section by construction. The debt share of AI capital spending
is taken from an economy-wide stock ratio, \result{DebtShareLow} to \result{DebtShareHigh} percent,
standing in for a marginal flow; the evidence in Section~\ref{sec:bust} points below that range,
and because the rate on the debt-financed normal return is negative, a lower debt share would raise
the assembled rate and harden the finding rather than soften it. And the condition is about
revenue, not about welfare or about whether any of this should be taxed more. What follows from a
shortfall is a question for Section~\ref{sec:institutions}.
